% Self-contained source. No images, external inputs, or external bibliography.
% Build from reports/: pdflatex -synctex=1 -interaction=nonstopmode -halt-on-error current_status_report.tex
% Run pdflatex twice to resolve the contents and cross-references. No BibTeX is required.
\documentclass[11pt,a4paper]{article}
\usepackage[T1]{fontenc}
\usepackage[utf8]{inputenc}
\usepackage{lmodern,microtype}
\usepackage[margin=25mm]{geometry}
\usepackage{amsmath,amssymb,amsthm,mathtools}
\usepackage{booktabs,tabularx,array}
\usepackage[hidelinks]{hyperref}
\setlength{\parskip}{4pt}
\setlength{\parindent}{0pt}
\setlength{\emergencystretch}{2em}
\newcommand{\ket}[1]{\lvert #1\rangle}
\newcommand{\bra}[1]{\langle #1\rvert}
\newcommand{\Tr}{\operatorname{Tr}}
\newcommand{\id}{\mathrm{id}}
\newcommand{\diag}{\operatorname{diag}}
\newcommand{\rank}{\operatorname{rank}}
\newcommand{\spec}{\operatorname{spec}}
\newcommand{\HH}{\mathcal H}
\newcommand{\ZZ}{\mathrm{ZZ}}
\newcommand{\ex}{\mathrm{ex}}
\newcommand{\OmF}{\Omega_{T_F}}
\newcommand{\OmR}{\Omega_R}
\newtheorem{proposition}{Proposition}[section]
\title{Two spins coupled to an XXZ bulk\\[3pt]
\large Reduced dynamics and quantum-state texture}
\author{Spin transport project}
\date{3 November 2026}

\begin{document}
\maketitle
\begin{abstract}
We study the joint reduced dynamics of two spin-$1/2$ probes,
each attached to one end of a finite XXZ bulk. The total Hamiltonian generates
unitary evolution, and the initial probe--bulk state is a product. For Ising
boundary coupling, four conditional bulk Hamiltonians give an exact
time-dependent Schur, or generalized dephasing, channel. Its complete action
is encoded in the output from the single probe preparation $\ket{++}$.
For exchange coupling, the available description gives the symmetry
blocks rather than closed expressions for their coefficients. When the Hamiltonian conserves total
magnetization and the bulk preparation commutes with bulk magnetization,
the reduced superoperator and Choi matrix split into blocks of sizes
$1,4,6,4,1$. Quantum-state texture is evaluated only on the output from
$\ket{++}$; exact relations to the Ising channel, its fidelities and its
entanglement-generating ability follow from the Schur structure.
\end{abstract}
\tableofcontents

\section{Introduction}

Two probes coupled to a common spin chain form an open quantum system
whose reduced dynamics depends on the interaction at each contact.
We consider Ising and exchange contacts separately, with no direct
probe--probe interaction. The bulk starts in a fixed state, uncorrelated
with the probe pair. Tracing out the bulk after unitary evolution defines
a quantum channel on the pair.

Ising contacts conserve both probe $z$ projections and produce generalized
dephasing in the joint computational basis. Exchange contacts allow spin
transfer between the probes and the bulk. Conservation of total
magnetization then constrains the reduced channel when the bulk
preparation respects the same rotation symmetry.

These two interactions lead to the following channel descriptions:
\begin{enumerate}
\item Ising coupling: an exact conditional-Hamiltonian reduction, an exact
Schur channel, and reconstruction of that channel from the full output
state for $\ket{++}$.
\item Exchange coupling: exact total-magnetization sectors and, for a
magnetization-invariant bulk preparation, a phase-covariant channel with
five coherence-order blocks. These constraints leave the block entries
to be computed from the many-body evolution.
\item Texture: exact grand-sum, return-probability and response-operator
identities. For Ising coupling these yield channel fidelities and other
properties without additional input preparations.
\end{enumerate}

The derivations use exact finite-system dynamics, without a weak-coupling
or Markov approximation. The channel structure fixes which probe
operators can mix; their time dependence depends on the bulk Hamiltonian
and its initial state.

\section{Model, conventions and the reduced channel}
\label{sec:model}

\subsection{Geometry and Hamiltonians}

Let $N\ge2$ be the number of \emph{bulk} spins. There are $N+2$ spins in
total, with physical order $A,1,\ldots,N,B$. We use
\begin{equation}
 S^\alpha=\sigma^\alpha/2,\qquad
 S^\pm=S^x\pm iS^y=(\sigma^x\pm i\sigma^y)/2,\qquad \hbar=1.
\end{equation}
The basis convention is $\ket0=\ket\uparrow$, $\ket1=\ket\downarrow$,
and the probe basis is $(\ket{00},\ket{01},\ket{10},\ket{11})$.
For formulas involving a product of probe and bulk operators, we group
the Hilbert space as $\HH_{AB}\otimes\HH_E$ by a fixed permutation
of the physical site ordering.

The reference bulk and the two alternative interactions are
\begin{align}
 H_E&=J\sum_{j=1}^{N-1}
 (S_j^xS_{j+1}^x+S_j^yS_{j+1}^y+\Delta S_j^zS_{j+1}^z),
 \label{eq:bulk}\\
 V_{\ZZ}&=g_A S_A^zS_1^z+g_B S_N^zS_B^z,
 \label{eq:ising}\\
 V_{\ex}&=g_A(S_A^xS_1^x+S_A^yS_1^y)
          +g_B(S_N^xS_B^x+S_N^yS_B^y)\notag\\
 &=\frac{g_A}{2}(S_A^+S_1^-+S_A^-S_1^+)
   +\frac{g_B}{2}(S_N^+S_B^-+S_N^-S_B^+).
 \label{eq:exchange}
\end{align}
Thus $H^{(r)}=H_E+V_r$, with $r=\ZZ$ or $\ex$.
Identity operators on the remaining spins are implicit. The parameters
$J,g_A,g_B$ have energy units and $\Delta$ is dimensionless. Equal
couplings mean $g_A=g_B=g$. The exchange flip-flop matrix element at
the left contact is $g_A/2$ in this convention.

Equation~\eqref{eq:bulk} has no applied fields. Longitudinal bulk or probe
fields preserve the magnetization arguments below, with their phases
included in the conditional dynamics. Transverse \emph{bulk} fields also
preserve the Ising Schur reduction, but generally invalidate the exchange
magnetization reduction because they break conservation of total
$z$ magnetization.

\subsection{Product preparation and exact channel}

Fix any bulk density matrix $\omega_E$, independent of the probe input.
The initial state and reduced evolution are
\begin{align}
 \rho(0)&=\rho_{AB}(0)\otimes\omega_E,\\
 U_t&=e^{-iH^{(r)}t},\qquad
 \Lambda_t(X)=\Tr_E[U_t(X\otimes\omega_E)U_t^\dagger],
 \label{eq:channel}\\
 \rho_{AB}(t)&=\Lambda_t(\rho_{AB}(0)).
\end{align}
The product assumption separates the pair from the bulk; it imposes no
factorization between $A$ and $B$ for the definition of the channel.
Only the texture protocol fixes their input to $\ket{++}$.

To see complete positivity and trace preservation directly, write
$\omega_E=\sum_\nu p_\nu\ket{\phi_\nu}\bra{\phi_\nu}$ and choose any
orthonormal bulk output basis $\{\ket r\}$. Then
\begin{equation}
 K_{r\nu}(t)=\sqrt{p_\nu}\,{}_E\bra r U_t\ket{\phi_\nu}_E,
 \quad \Lambda_t(X)=\sum_{r,\nu}K_{r\nu}XK_{r\nu}^\dagger,
 \quad \sum_{r,\nu}K_{r\nu}^\dagger K_{r\nu}=I_4.
 \label{eq:kraus}
\end{equation}
These are formal exact expressions in both regimes; evaluating them for
an interacting bulk remains a many-body problem. Correlated probe--bulk
preparations have reduced trajectories, but do not in general define
the same CPTP map on arbitrary probe inputs.

\subsection{Choi normalization}

The \emph{normalized} Choi state, with the output factor first, is
\begin{equation}
 \ket{\Phi_4}=\frac12\sum_{a=1}^4\ket a_{AB}\ket a_R,
 \qquad J_t=(\Lambda_t\otimes\id_R)(\ket{\Phi_4}\bra{\Phi_4}).
 \label{eq:choi}
\end{equation}
It obeys $J_t\succeq0$, $\Tr J_t=1$ and
$\Tr_{AB}J_t=I_R/4$. The corresponding unnormalized Choi matrix is
$C_t=4J_t$. The auxiliary system $R$ represents a reference for the
channel and is independent of the physical bulk.

\section{Ising coupling: exact generalized dephasing}
\label{sec:zz}

\subsection{Four conditional Hamiltonians}

Since $[H^{(\ZZ)},S_A^z]=[H^{(\ZZ)},S_B^z]=0$, label each probe basis
state by $a=(s_A,s_B)$ with $s_A,s_B=\pm1/2$, and set
$P_a=\ket a\bra a$. The Hamiltonian and propagator are
\begin{align}
 H^{(\ZZ)}&=\sum_a P_a\otimes H_a,
 &H_a&=H_E+g_A s_A S_1^z+g_B s_B S_N^z,\label{eq:conditional}\\
 U_t&=\sum_aP_a\otimes U_a(t),
 &U_a(t)&=e^{-iH_at}.
\end{align}
This follows by inserting the spectral decompositions of $S_A^z,S_B^z$;
orthogonal projectors exponentiate independently. It is exact for any
bulk Hamiltonian acting only on $E$, including one with transverse fields.

For equal couplings, the four perturbations are
\begin{equation}
\begin{array}{c|c}
 a&H_a-H_E\\\hline
 00&\frac g2(S_1^z+S_N^z)\\
 01&\frac g2(S_1^z-S_N^z)\\
 10&-\frac g2(S_1^z-S_N^z)\\
 11&-\frac g2(S_1^z+S_N^z)
\end{array}.
\end{equation}
These configurations label invariant branches. The input may still
coherently occupy all four branches.

\subsection{Schur form and its physical constraints}

Let $E_{ab}=\ket a\bra b$. Since
$P_cE_{ab}P_d=\delta_{ca}\delta_{db}E_{ab}$, tracing the controlled
propagator gives
\begin{equation}
 \boxed{\Lambda_t(E_{ab})=\Gamma_{ab}(t)E_{ab},\qquad
 \Lambda_t(X)=\Gamma(t)\circ X,}
 \label{eq:schur}
\end{equation}
where $\circ$ is the entrywise Schur product and
\begin{equation}
 \Gamma_{ab}(t)=\Tr_E[U_a(t)\omega_E U_b(t)^\dagger].
 \label{eq:gamma}
\end{equation}
For a pure bulk, this is the complex conditional overlap
$\bra{\phi_0}U_b^\dagger U_a\ket{\phi_0}$. Its squared modulus is a
conditional Loschmidt echo; the complex amplitude, including its phase,
is needed for the channel.

The matrix $\Gamma$ is a correlation matrix:
\begin{equation}
 \Gamma=\Gamma^\dagger\succeq0,\qquad \Gamma_{aa}=1,
 \qquad |\Gamma_{ab}|\le1,\qquad \Gamma_{ab}(0)=1.
 \label{eq:correlation}
\end{equation}
Indeed, setting $k_{r\nu,a}=\sqrt{p_\nu}\bra rU_a\ket{\phi_\nu}$
gives $\Gamma_{ab}=\sum_{r,\nu}k_{r\nu,a}k_{r\nu,b}^*$, a positive
Gram factorization. The unit diagonal follows from unitarity and
$\Tr\omega_E=1$; $2\times2$ principal minors give the modulus bound.
Conversely, every positive matrix with unit diagonal defines a CPTP
Schur map. This converse characterizes the map class, not which members
a specified XXZ Hamiltonian realizes at a given time.

All four populations are fixed and $\Lambda_t(I_4)=I_4$. Six complex
off-diagonal factors specify the nontrivial action, subject to positivity.
The common interacting bulk can correlate the two probes, so a product
of independent single-probe dephasing channels cannot be assumed.
There is no general monotonicity of $|\Gamma_{ab}(t)|$.

\subsection{Dephasing and depolarization}\label{sec:depolarizing}

A depolarizing map has the form
$\mathcal D_\lambda(X)=\lambda X+(1-\lambda)\Tr(X)I_4/4$.
The Ising map fixes every $P_a$. Requiring
$\mathcal D_\lambda(P_a)=P_a$ forces $\lambda=1$.
Thus a nonidentity Ising Schur map cannot be a depolarizing map.
Dephasing includes coherent diagonal rotations and revivals; it does
not imply irreversible or strictly dissipative evolution.

\subsection{One output determines the complete channel}
\label{sec:oneinput}

The sole probe preparation used for texture is
\begin{equation}
 \ket{f_1}=\ket+_A\otimes\ket+_B
 =\frac12\sum_{a=1}^4\ket a,\qquad
 f_1=\ket{f_1}\bra{f_1},\qquad (f_1)_{ab}=\frac14.
 \label{eq:f1}
\end{equation}
Denote its full output state by $\rho_f(t)=\Lambda_t(f_1)$. Then
\begin{equation}
 \boxed{\rho_f(t)=\Gamma(t)/4,\qquad
 \Gamma(t)=4\rho_f(t),\qquad
 \Lambda_t(X)=[4\rho_f(t)]\circ X.}
 \label{eq:reconstruction}
\end{equation}
Consequently the \emph{entire density matrix} of one known output
determines every property of the instantaneous Ising channel. Temporal
properties require these outputs at the relevant times. A single scalar
texture value cannot replace this matrix.

The populations of $\rho_f$ are $1/4$, so its $12$ remaining real
parameters are its six complex coherences. In a Pauli expansion,
$\langle ZI\rangle=\langle IZ\rangle=\langle ZZ\rangle=0$.
The eight local measurement settings $(\alpha,\beta)\ne(z,z)$,
$\alpha,\beta\in\{x,y,z\}$, recover the other Pauli coefficients,
including single-spin expectations from their marginals. Thus one
input preparation together with output tomography suffices.

\subsection{Choi, Kraus and spectral consequences}

Let $W:\HH_{AB}\to\HH_{AB}\otimes\HH_R$ be the isometry
$W\ket a=\ket{aa}$. Equations~\eqref{eq:choi} and \eqref{eq:schur} give
\begin{equation}
 J_t=\frac14\sum_{a,b}\Gamma_{ab}\ket{aa}\bra{bb}
     =W\rho_f(t)W^\dagger.
 \label{eq:choizz}
\end{equation}
The Choi support has dimension at most four. In particular,
\begin{align}
 \rank J_t&=\rank\Gamma\le4,\\
 \spec(J_t)&=\spec(\Gamma/4)\ \text{with twelve additional zeros},\\
 S(J_t)&=S(\rho_f),\qquad
 \Tr J_t^2=\Tr\rho_f^2=\frac1{16}\sum_{a,b}|\Gamma_{ab}|^2.
 \label{eq:choipurity}
\end{align}
If $\Gamma=\sum_k\lambda_kv_kv_k^\dagger$ with normalized eigenvectors,
one minimal Kraus family is
\begin{equation}
 D_k=\sqrt{\lambda_k}\,\diag(v_{k,1},\ldots,v_{k,4}),
 \qquad \Lambda_t(X)=\sum_k D_kXD_k^\dagger.
\end{equation}
The superoperator is diagonal in the $E_{ab}$ basis, with eigenvalues
$\Gamma_{ab}$ and singular values $|\Gamma_{ab}|$. Choi rank and
superoperator rank are distinct: the identity has Choi rank one and
superoperator rank sixteen.

The map is unitary iff $\rank\Gamma=1$, in which case
$\Gamma_{ab}=u_au_b^*$ with $|u_a|=1$, giving the diagonal unitary
$\diag(u_1,\ldots,u_4)$. This does not require that the unitary factor
into two local rotations.

The partial transpose of $J_t$ has eigenvalues $1/4$ (four times) and
$\pm|\Gamma_{ab}|/4$ for each $a<b$. Hence its reference--output
negativity is
\begin{equation}
 \mathcal N_{AB|R}(J_t)=\frac14\sum_{a<b}|\Gamma_{ab}|.
\end{equation}
This is distinct from physical $A|B$ entanglement. The Ising map is
entanglement breaking iff all off-diagonal $\Gamma_{ab}$ vanish:
otherwise its Choi state has a negative partial transpose, while the
diagonal case has the explicitly separable Choi state
$\frac14\sum_a\ket{aa}\bra{aa}$.

\subsection{Time-local dynamics and intermediate maps}

The map is invertible iff every $\Gamma_{ab}(t)$ is nonzero. At regular
times its exact time-local generator obeys
\begin{equation}
 \mathcal L_t(E_{ab})=\frac{\dot\Gamma_{ab}(t)}{\Gamma_{ab}(t)}E_{ab}.
\end{equation}
If $\Lambda_s$ is invertible, the unique intermediate linear map is
Schur multiplication by $Q_{ab}(t,s)=\Gamma_{ab}(t)/\Gamma_{ab}(s)$.
Its diagonal is one, so
\begin{equation}
 \Lambda_t\Lambda_s^{-1}\text{ is CPTP}\quad\Longleftrightarrow\quad
 Q(t,s)\succeq0.
 \label{eq:divisibility}
\end{equation}
At singular times the division is undefined and this test cannot be
applied. A vanished coherence factor that later becomes nonzero already
precludes divisibility through that time, since a linear intermediate
map cannot recover the annihilated matrix unit.

\subsection{Probe entanglement from the same input}
\label{sec:entangling}

For two-qubit states, concurrence $C$ and negativity $\mathcal N$ test
entanglement of $A$ with $B$; their definitions are in
Appendix~\ref{app:observables}. For a Schur channel, the output from
$\ket{++}$ also determines the maximum concurrence generated from
separable probe states.

\begin{proposition}
For a fixed Ising Schur channel, its maximum output concurrence over
separable probe inputs is
\begin{equation}
 \max_{\rho\in\mathrm{Sep}(A:B)}C[\Lambda_t(\rho)]
 =C[\rho_f(t)]=C[\Gamma(t)/4].
 \label{eq:cgen}
\end{equation}
\end{proposition}
\begin{proof}
For a normalized pure product input
$(a_0\ket0+a_1\ket1)\otimes(b_0\ket0+b_1\ket1)$, set
$A=\sqrt2\diag(a_0,a_1)$ and $B=\sqrt2\diag(b_0,b_1)$.
The output is $(A\otimes B)\rho_f(A\otimes B)^\dagger$ and has unit
trace because $\rho_f$ has uniform diagonal. The local-filter identity
for two-qubit concurrence gives
$C_{\mathrm{out}}=|\det A\det B|C(\rho_f)
=4|a_0a_1b_0b_1|C(\rho_f)\le C(\rho_f)$.
The identity follows from the spin-flip formula; singular filters follow
by continuity. Convexity extends the bound to mixtures of product
inputs. The input $\ket{++}$ attains equality.
\end{proof}
In particular, the Ising map can create probe entanglement from some
separable input iff $\rho_f(t)$ is entangled. Texture and Choi
reference--output negativity are different quantities and do not replace
this criterion.

\section{Exchange coupling: exact symmetry reduction}
\label{sec:ex}

\subsection{Total charge and Hamiltonian sectors}

For the field-free model define
\begin{equation}
 M_S=S_A^z+S_B^z,\qquad M_E=\sum_{j=1}^N S_j^z,
 \qquad M=M_S+M_E.
\end{equation}
Each exchange term transfers one unit between a probe and the bulk,
so $[H^{(\ex)},M]=0$, while individual probe $S^z$ values generally
are not conserved. The total-charge sector $m$ has dimension
\begin{equation}
 \dim\HH_m=\binom{N+2}{(N+2)/2-m}.
\end{equation}
This Hamiltonian decomposition requires no assumption on the input;
coherences between occupied sectors must still be retained.
There is no four-branch controlled propagator of the form
\eqref{eq:conditional} for exchange coupling.

\subsection{Bulk assumption and proof of phase covariance}

For the channel reduction assume
\begin{equation}
 [\omega_E,M_E]=0.
 \label{eq:bulksym}
\end{equation}
Support in a single well-defined bulk magnetization sector is sufficient.
An incoherent mixture of sectors also suffices. A coherent superposition
of different bulk magnetizations generally does not.

Let $V_\theta=e^{-i\theta M_S}$ and $W_\theta=e^{-i\theta M_E}$.
Conservation gives $[U_t,V_\theta\otimes W_\theta]=0$, while
\eqref{eq:bulksym} gives $W_\theta\omega_EW_\theta^\dagger=\omega_E$.
Inserting these into the partial trace proves
\begin{equation}
 \boxed{\Lambda_t(V_\theta XV_\theta^\dagger)
       =V_\theta\Lambda_t(X)V_\theta^\dagger.}
 \label{eq:covariance}
\end{equation}
This is $U(1)$ phase covariance. It permits exchange of magnetization
with the bulk; it does not imply $\Lambda_t^*(M_S)=M_S$ or unitality.

\subsection{Five superoperator blocks}

The probe charges are $(m_{00},m_{01},m_{10},m_{11})=(1,0,0,-1)$.
The operator $E_{ab}$ has coherence order $q=m_a-m_b$, meaning
$[M_S,E_{ab}]=qE_{ab}$. Under a rotation it acquires the phase
$e^{-iq\theta}$, so covariance implies
\begin{equation}
 \Lambda_t(\mathcal V_q)\subseteq\mathcal V_q,\qquad
 \mathcal V_q=\operatorname{span}\{E_{ab}:m_a-m_b=q\}.
\end{equation}
The ordered bases $\mathcal B_q$ of these spaces are
\begin{center}
\begin{tabular}{@{}ccl@{}}
\toprule
$q$&Dimension&Basis of $\mathcal V_q$\\\midrule
$0$&6&$E_{00,00},E_{01,01},E_{10,10},E_{11,11},E_{01,10},E_{10,01}$\\
$+1$&4&$E_{00,01},E_{00,10},E_{01,11},E_{10,11}$\\
$-1$&4&Adjoints of the $q=+1$ basis\\
$+2$&1&$E_{00,11}$\\
$-2$&1&$E_{11,00}$\\\bottomrule
\end{tabular}
\end{center}
Thus the $16\times16$ superoperator has the exact reduction
\begin{equation}
 \boxed{\Lambda_t\simeq T^{(-2)}\oplus T^{(-1)}\oplus T^{(0)}
               \oplus T^{(+1)}\oplus T^{(+2)},\quad
               (d_q)=(1,4,6,4,1).}
 \label{eq:blocks}
\end{equation}
Operators of the same order may mix. Hermiticity relates the opposite
orders. In the ordered basis $\mathcal B_q$, the exact entries are
\begin{equation}
 T^{(q)}_{cd,ab}(t)
 =\sum_{r,\nu}(K_{r\nu})_{ca}(K_{r\nu})_{db}^*,
 \qquad E_{ab},E_{cd}\in\mathcal B_q.
 \label{eq:blockentries}
\end{equation}
Equation~\eqref{eq:blockentries} is a formal construction, not a closed
evaluation of these functions for general interacting XXZ dynamics.

The normalized Choi state satisfies the equivalent symmetry condition
\begin{equation}
 [J_t,M_S\otimes I_R-I_{AB}\otimes M_S^{\mathsf T}]=0.
\end{equation}
Grouping output--reference kets $\ket c\ket a$ by $m_c-m_a$ gives Choi
blocks of the same sizes $1,4,6,4,1$. Choi blocks and superoperator
blocks are different rearrangements of channel coefficients and should
not be identified as matrices or assigned the same spectra.

\subsection{Kraus operators with definite charge transfer}

Under \eqref{eq:bulksym}, diagonalize $\omega_E$ within its charge
sectors: $M_E\ket{\phi_\nu}=\mu_\nu\ket{\phi_\nu}$. Choose
$M_E\ket r=\mu_r\ket r$. Conservation of the total input and output
charge in Eq.~\eqref{eq:kraus} yields
\begin{equation}
 (K_{r\nu})_{ca}\ne0\ \Longrightarrow\
 m_c-m_a=\mu_\nu-\mu_r,\qquad
 [M_S,K_{r\nu}]=(\mu_\nu-\mu_r)K_{r\nu}.
\end{equation}
Only transfers between $-2$ and $+2$ can act on the two probes.
The transfer label belongs to a Kraus operator; coherence order labels
an input or output matrix unit. Their roles in the Choi and
superoperator reductions are related but distinct.

\subsection{Reduced states at fixed total magnetization}
\label{sec:xstate}

If the \emph{global} initial state is supported in one total-$M$ sector,
the reduced probe state has the exact form
\begin{equation}
 \rho_{AB}(t)=\begin{pmatrix}
 p_{00}&0&0&0\\0&p_{01}&c&0\\0&c^*&p_{10}&0\\0&0&0&p_{11}
 \end{pmatrix}.
 \label{eq:xstate}
\end{equation}
For a pure global ket, resolve its four probe components. Their bulk
kets have charges $m-1,m,m,m+1$. Orthogonality of different bulk charges
eliminates every coherence except $01\leftrightarrow10$. Linearity
extends the conclusion to mixtures within the sector.
The populations and $c$ remain time dependent, and the concurrence is
\begin{equation}
 C=2\max\{0,|c|-\sqrt{p_{00}p_{11}}\}.
\end{equation}

This restriction does \emph{not} apply to the texture protocol merely
because the bulk has fixed magnetization. The input $\ket{++}$ contains
probe charges $+1,0,-1$ with probabilities $1/4,1/2,1/4$, and hence
occupies three total-charge sectors for a fixed-charge bulk. Their
coherences cannot be discarded. In particular, $f_1$ itself is not
of the form~\eqref{eq:xstate}.

Reflection and global spin flip give additional covariance only when
both the Hamiltonian and bulk preparation have the required symmetry.
Equal contacts or a symmetric Hamiltonian alone do not impose that
symmetry on an arbitrary reduced trajectory.

\section{Quantum-state texture of the \texorpdfstring{$\ket{++}$}{++} preparation}
\label{sec:texture}

\subsection{Grand sum and directly measurable return probability}

Texture is defined relative to the probe computational basis, whose
textureless state is $f_1$ in Eq.~\eqref{eq:f1}. For any probe state,
the grand sum is real and satisfies
\begin{equation}
 \Sigma(\rho)=\sum_{a,b}\rho_{ab}=4\bra{f_1}\rho\ket{f_1},
 \qquad 0\le\Sigma(\rho)\le4.
\end{equation}
For the single physical input $f_1\otimes\omega_E$, set
\begin{equation}
 p_f(t)=\bra{f_1}\rho_f(t)\ket{f_1}=\Sigma(\rho_f(t))/4,
 \quad T_F(\rho_f)=1-p_f,\quad R(\rho_f)=-\ln p_f.
 \label{eq:textures}
\end{equation}
These are the fidelity texture and rugosity conventions used here
\cite{Parisio2024,Vieira2026}. At $p_f=0$, $R=+\infty$; initially
$p_f=1$ and both textures vanish. The texturizing powers evaluated on
the unique textureless input are $\OmF(t)=T_F[\Lambda_t(f_1)]$ and
$\OmR(t)=R[\Lambda_t(f_1)]$.

Since $4f_1=(I+\sigma_A^x)(I+\sigma_B^x)$,
\begin{equation}
 \Sigma(\rho_f)=1+\langle\sigma_A^x\rangle
 +\langle\sigma_B^x\rangle+\langle\sigma_A^x\sigma_B^x\rangle.
 \label{eq:measurement}
\end{equation}
A joint $x$-basis measurement gives $p_f=P_{++}^{(x)}$ directly.
Equivalently, apply a Hadamard to each probe and record outcome $00$.
This uses one measurement setting, whereas full channel reconstruction
in the Ising case uses the additional coherence measurements described
in Section~\ref{sec:oneinput}.

\subsection{Ising texture and channel fidelities}

Using $\rho_f=\Gamma/4$ gives
\begin{align}
 p_f&=\frac1{16}\sum_{a,b}\Gamma_{ab}
     =\frac14+\frac18\sum_{a<b}\Re\Gamma_{ab},\\
 \OmF&=\frac34-\frac18\sum_{a<b}\Re\Gamma_{ab},
 &\OmR&=-\ln\!\left(\frac1{16}\sum_{a,b}\Gamma_{ab}\right).
 \label{eq:omegazz}
\end{align}
For the same Schur channel, the entanglement fidelity relative to the
identity is $F_e=\bra{\Phi_4}J_t\ket{\Phi_4}=p_f$. Therefore the
average pure-input fidelity relative to the identity is
\begin{equation}
 \boxed{F_e=p_f=\Sigma(\rho_f)/4,\qquad
 F_{\mathrm{avg}}=\frac{4F_e+1}{5}
                 =\frac{1+\Sigma(\rho_f)}5.}
 \label{eq:fidelities}
\end{equation}
The general relation between average and entanglement fidelities is
standard \cite{Nielsen2002}; $F_e=p_f$ uses the Schur structure.
These are fidelities of the \emph{two-probe map relative to the identity},
not fidelities of sending an unknown state from $A$ to $B$.

Only the real sum of the six coherences enters $p_f$. It does not fix
their individual magnitudes or phases and does not determine purity,
probe entanglement, or CP divisibility. Coherent phase evolution can
change texture even when the map is unitary.

\subsection{Response operator and texture preservation}

The dual channel gives
\begin{equation}
 \mathcal R_t=4\Lambda_t^*(f_1),\qquad
 \Sigma[\Lambda_t(X)]=\Tr(\mathcal R_tX).
 \label{eq:response}
\end{equation}
In the Ising case, comparing coefficients of $X_{ab}$ yields
\begin{equation}
 \mathcal R_t=\Gamma(t)^{\mathsf T}.
 \label{eq:responsetranspose}
\end{equation}
The transpose is ordinary, not a Hermitian conjugate. In general,
fixing the free state, $\Lambda_t(f_1)=f_1$, and preserving the grand
sum for every input, $\Lambda_t^*(f_1)=f_1$, are different conditions
\cite{Vieira2026}. For this Schur family each condition forces
$\Gamma_{ab}=1$ for all $a,b$, hence the identity channel. Consequently
$\OmF=0$ also certifies the identity within the known Ising family.

\subsection{Exchange texture and the limits of one input}

Equations~\eqref{eq:textures}, \eqref{eq:measurement} and
\eqref{eq:response} hold for exchange coupling as well. However,
$\rho_f$ is generally not a Schur multiplier matrix and does not
determine the exchange channel. Under Eq.~\eqref{eq:bulksym}, $f_1$
has a nonzero component in every $\mathcal V_q$, but one output
determines only each block's action on that particular component.
It cannot determine a general $6\times6$ or $4\times4$ block.

There is one exact scalar exception. Since $\mathcal V_{+2}$ is
one-dimensional, write
$\Lambda_t(E_{00,11})=\eta_{+2}(t)E_{00,11}$. Then
\begin{equation}
 \eta_{+2}(t)=4[\rho_f(t)]_{00,11},\qquad
 \eta_{-2}(t)=\eta_{+2}(t)^*.
\end{equation}
This inference requires covariance; it is not guaranteed for a bulk
coherent across magnetization sectors or for transverse-field dynamics.
The equality $F_e=p_f$ is likewise not a general exchange-channel
identity, although $F_{\mathrm{avg}}=(4F_e+1)/5$ still holds.

\section{Channel structure and dynamical coefficients}

Six complex overlaps between conditional bulk evolutions determine the
Ising channel. The output $\rho_f=\Gamma/4$ therefore specifies its full
action, including its Choi state, fidelities and ability to generate
probe entanglement. Evaluating the overlaps requires the conditional
dynamics for the chosen $J,\Delta,g_A,g_B,N$ and $\omega_E$.

For exchange coupling, conservation of total magnetization and a
magnetization-invariant bulk preparation give five coherence-order
blocks. Symmetry fixes their dimensions and forbids mixing between
different orders, but leaves the entries within each block to be
calculated. The output from $\ket{++}$ determines the texture and the
one-dimensional highest-order response; it does not specify the full
exchange channel. A general closed expression for the interacting
exchange coefficients remains unavailable.

\appendix
\section{State diagnostics and their interpretation}\label{app:observables}

For $\rho=\rho_{AB}$, the state purity and entropy are
\begin{equation}
 P=\Tr\rho^2,\qquad S(\rho)=-\sum_k\lambda_k\ln\lambda_k,
 \qquad 0\ln0=0.
\end{equation}
Logarithms are natural. The probe mutual information is
$I(A:B)=S(\rho_A)+S(\rho_B)-S(\rho)$, with
$\rho_A=\Tr_B\rho$ and $\rho_B=\Tr_A\rho$. It includes classical
and quantum correlations.

Probe negativity and concurrence are
\begin{align}
 \mathcal N_{A|B}(\rho)&=\frac{\|\rho^{T_B}\|_1-1}{2},\\
 \widetilde\rho&=(\sigma_y\otimes\sigma_y)\rho^*
 (\sigma_y\otimes\sigma_y),\qquad
 C(\rho)=\max(0,r_1-r_2-r_3-r_4),
\end{align}
where the $r_i$ are the decreasing nonnegative square roots of the
eigenvalues of $\rho\widetilde\rho$ \cite{Wootters1998}.
For two qubits, $C>0$ and $\mathcal N>0$ are equivalent tests of
entanglement. They concern $A|B$, while $S(\rho_{AB})$ concerns the
pair's marginal mixedness. Only for a globally pure state does
$S(\rho_{AB})$ equal entanglement entropy across $AB|E$.

For the Ising texture input,
\begin{equation}
 P_f(t)=\frac14+\frac18\sum_{a<b}|\Gamma_{ab}(t)|^2,
 \qquad \frac14\le P_f\le1.
\end{equation}
Its local probe coherences are
\begin{equation}
 (\rho_A)_{01}=\frac{\Gamma_{00,10}+\Gamma_{01,11}}4,
 \qquad (\rho_B)_{01}=\frac{\Gamma_{00,01}+\Gamma_{10,11}}4,
\end{equation}
and each local diagonal is $(1/2,1/2)$.
For a pure bulk, the global state is
$\frac12\sum_a\ket a\otimes U_a\ket{\phi_0}$; it is a product
across $AB|E$ iff its four conditional bulk kets are proportional.
For a mixed bulk, reduced entropy or purity alone does not quantify
entanglement across that cut.

\begin{thebibliography}{9}
\bibitem{Parisio2024}
F. Parisio, ``Quantum-State Texture and Gate Identification,''
\emph{Phys. Rev. Lett.} \textbf{133}, 260801 (2024).
\href{https://doi.org/10.1103/PhysRevLett.133.260801}{doi:10.1103/PhysRevLett.133.260801}.
\bibitem{Vieira2026}
C. H. S. Vieira, X. Nie, D. Lu and F. Parisio,
``Quantum-State Texture Dynamics: Theory and Experiment,''
\href{https://arxiv.org/abs/2609.05248}{arXiv:2609.05248} (2026).
\bibitem{Nielsen2002}
M. A. Nielsen, ``A simple formula for the average gate fidelity of a
quantum dynamical operation,'' \emph{Phys. Lett. A} \textbf{303},
249--252 (2002).
\href{https://arxiv.org/abs/quant-ph/0205035}{arXiv:quant-ph/0205035}.
\bibitem{Wootters1998}
W. K. Wootters, ``Entanglement of Formation of an Arbitrary State of Two
Qubits,'' \emph{Phys. Rev. Lett.} \textbf{80}, 2245--2248 (1998).
\href{https://doi.org/10.1103/PhysRevLett.80.2245}{doi:10.1103/PhysRevLett.80.2245}.
\end{thebibliography}
\end{document}
